\documentclass[11pt,a4paper]{ctexart}

\usepackage[margin=2.2cm]{geometry}
\usepackage{amsmath,amssymb,amsthm,mathtools}
\usepackage[dvipsnames]{xcolor}
\usepackage{hyperref}
\usepackage{enumitem}
\usepackage[most]{tcolorbox}
\usepackage{tikz}
\usepackage{listings}
\usepackage{array}
\usepackage{booktabs}
\usetikzlibrary{arrows.meta,calc,matrix,positioning}

\hypersetup{
  colorlinks=true,
  linkcolor=blue!55!black,
  urlcolor=blue!55!black
}

\setlist[enumerate]{itemsep=0.35em, topsep=0.45em}
\setlist[itemize]{itemsep=0.35em, topsep=0.45em}

\newtheorem{theorem}{定理}[section]
\newtheorem{definition}[theorem]{定义}
\newtheorem{example}[theorem]{例}
\newtheorem{remark}[theorem]{备注}

\tcbset{
  commonbox/.style={
    enhanced,
    breakable,
    boxrule=0.8pt,
    arc=2mm,
    left=1.4mm,
    right=1.4mm,
    top=1mm,
    bottom=1mm,
    colback=white
  }
}

\newtcolorbox{problemblock}{commonbox,colframe=BrickRed!65!black,colback=BrickRed!4,title=问题,fonttitle=\bfseries}
\newtcolorbox{intuitionblock}{commonbox,colframe=RoyalBlue!70!black,colback=RoyalBlue!4,title=朴素直觉,fonttitle=\bfseries}
\newtcolorbox{mathbox}[1]{commonbox,colframe=Purple!60!black,colback=Purple!4,title=#1,fonttitle=\bfseries}
\newtcolorbox{codebox}{commonbox,colframe=ForestGreen!60!black,colback=ForestGreen!4,title=代码对应,fonttitle=\bfseries}
\newtcolorbox{keypointblock}{commonbox,colframe=black!65,colback=black!3,title=记住这一点,fonttitle=\bfseries}

\lstset{
  basicstyle=\ttfamily\small,
  breaklines=true,
  columns=fullflexible,
  keywordstyle=\color{Purple!70!black}\bfseries,
  commentstyle=\color{gray!70!black}\itshape,
  stringstyle=\color{ForestGreen!50!black},
  showstringspaces=false,
  language=Python,
  frame=single,
  rulecolor=\color{black!20},
  framesep=2pt,
  xleftmargin=0.5em
}

\newcommand{\vect}[1]{\boldsymbol{#1}}

\title{零基础读懂 \texttt{ch03.ipynb}\\\large 从「拿一句话喂给注意力」到「多头注意力」}
\author{}
\date{\today}

\begin{document}

\maketitle
\tableofcontents
\newpage

\section{这一章在解决什么}

\begin{problemblock}
GPT 一类的语言模型在做一件事：给定前面已经写好的若干个词，
预测下一个词是什么。要让模型「猜得准」，每个位置都需要参考它前面所有的词。
最早的做法是用循环网络（RNN）一个一个 token 处理，最后挤出一个固定大小的隐藏状态。
当句子变长，这个隐藏状态就变成了瓶颈：远处的信息要么被压扁，要么被忘掉。

\textbf{注意力机制（Attention）就是为了解决这个瓶颈而生：}
让每个位置可以\emph{直接}看其它每个位置，根据「相关程度」加权，再把信息汇总过来。

第 3 章干的事，就是把「注意力」从最简单的雏形，一步一步搭成 GPT 真正使用的那个版本。
\end{problemblock}

\begin{intuitionblock}
你可以把一句话想成 6 个人坐成一排开会。每个人都有一句话要说（输入向量），
但他在最终发言之前，会先\emph{扫一眼}屋里所有人，根据「这个人对我有多重要」打一个权重，
再把大家的话按权重加起来，作为他真正的「上下文版本」。
注意力机制就是把这个「扫一眼 + 加权汇总」用数学和代码写出来。
\end{intuitionblock}

\begin{keypointblock}
本章的最终交付物是一个叫 \texttt{MultiHeadAttention} 的 PyTorch 模块。
后面所有章节的 GPT 都直接拿它当积木。
\end{keypointblock}

\section{讲解原则}

为了不一上来就被符号淹没，本讲义对每一节都按下面四步展开：

\begin{enumerate}[label=\textbf{\arabic*.}]
  \item \textbf{背景问题}：上一节做完了，新的麻烦是什么？
  \item \textbf{朴素直觉}：先用人话说清楚，不引入新符号。
  \item \textbf{数学表达}：把直觉写成公式，每个新符号都立刻翻译成中文。
  \item \textbf{代码对应}：在 notebook 里这一段长什么样、张量形状是什么。
\end{enumerate}

如果某一节让你觉得「公式比直觉先到」，回到第 2 步。

\section{零基础最该问的 7 个问题}

\begin{enumerate}[label=\textbf{Q\arabic*.}]
  \item 什么叫「一个词的向量表示」？为什么句子可以变成一堆数？
  \item 「注意力分数」到底在量什么？为什么是点积？
  \item 为什么要 softmax？不能直接归一化吗？
  \item Query、Key、Value 这三个词差在哪？为什么要分三份？
  \item 为什么后面要除以 $\sqrt{d_k}$？
  \item 「因果掩码（causal mask）」是干嘛的？为什么用 $-\infty$？
  \item 「多头」到底多在哪？是把网络做大了吗？
\end{enumerate}

下面所有节都在回答其中一个或几个问题。

\section{贯穿全章的玩具例子}

\begin{problemblock}
全章都用同一个 6 个词的英文小句作为玩具：

\begin{center}
\textit{Your\quad journey\quad starts\quad with\quad one\quad step}
\end{center}

每个词被替换成一个 3 维向量（这一步在第 2 章里讲过，叫 \emph{embedding}），
拼起来就是一个形状 $(6, 3)$ 的矩阵 \texttt{inputs}。
\end{problemblock}

\begin{codebox}
\begin{lstlisting}
inputs = torch.tensor(
  [[0.43, 0.15, 0.89],   # Your    x^(1)
   [0.55, 0.87, 0.66],   # journey x^(2)
   [0.57, 0.85, 0.64],   # starts  x^(3)
   [0.22, 0.58, 0.33],   # with    x^(4)
   [0.77, 0.25, 0.10],   # one     x^(5)
   [0.05, 0.80, 0.55]]   # step    x^(6)
)
\end{lstlisting}
\end{codebox}

\begin{intuitionblock}
\textbf{这些数字是什么意思？}\\
不要纠结具体数值——它们只是「每个词在 3 个隐含维度上的强度」。
真实模型里这些维度是几百到几千，并且是训练学出来的。
本章不学怎么得到这些数，本章学的是：\emph{已经有了这些数之后怎么用注意力机制把它们重新汇总}。
\end{intuitionblock}

之后凡是要看「中间结果长什么样」，几乎都用第 2 个词 \textit{journey}（也就是 \texttt{inputs[1]}）当主角。

\section{第一步：没有任何参数的注意力（§3.3）}

\subsection{背景问题}

我们想给每个词产出一个新的向量 $\vect{z}^{(i)}$，叫做\textbf{上下文向量（context vector）}。
它必须同时融合\emph{所有 6 个词}的信息，但不同词的贡献要不一样。
最朴素的问题是：「贡献的权重从哪来？」

\subsection{朴素直觉}

\begin{intuitionblock}
能不能直接用「两个词的相似度」当权重？\\
两个向量越像，它们的\emph{点积}就越大。所以一个最朴素的想法是：

\begin{enumerate}
  \item 拿当前这个词 $\vect{x}^{(i)}$ 跟所有 6 个词分别做点积，得到 6 个分数；
  \item 把 6 个分数压成「和为 1」的权重；
  \item 用权重把 6 个原始向量加起来。
\end{enumerate}

整个过程\emph{没有任何要学习的参数}，纯计算。
\end{intuitionblock}

\subsection{数学表达}

以 $\vect{x}^{(2)}$ 为「查询者」为例。

\begin{mathbox}{第一步：注意力分数（attention scores）}
\[
  \omega_{2i} \;=\; \vect{x}^{(i)} \cdot \vect{x}^{(2)},
  \qquad i = 1, 2, \dots, 6.
\]
\end{mathbox}

每个 $\omega_{2i}$ 是一个数：第 $i$ 个词跟第 2 个词有多「像」。

\begin{mathbox}{第二步：注意力权重（attention weights）}
\[
  \alpha_{2i} \;=\; \frac{\exp(\omega_{2i})}{\sum_{j=1}^{6} \exp(\omega_{2j})}.
\]
这就是 \textbf{softmax}：把任意实数压成「全为正、加起来等于 1」的概率。
\end{mathbox}

\begin{remark}
为什么不能直接做 $\alpha_{2i} = \omega_{2i} / \sum_j \omega_{2j}$？
因为分数可能是负数，相除以后会出现「负权重」，失去「权重」的语义。
softmax 先用 $\exp$ 把所有数变正，再归一化，所以无论原始分数多大多小，结果都是合法的概率分布。
\end{remark}

\begin{mathbox}{第三步：加权求和}
\[
  \vect{z}^{(2)} \;=\; \sum_{i=1}^{6} \alpha_{2i}\, \vect{x}^{(i)}.
\]
$\vect{z}^{(2)}$ 是一个新的 3 维向量，它就是「以 $\vect{x}^{(2)}$ 为视角看到的整句话」。
\end{mathbox}

\subsection{代码对应}

\begin{codebox}
\begin{lstlisting}
query = inputs[1]                       # x^(2)
attn_scores_2 = inputs @ query          # 6 个分数
attn_weights_2 = torch.softmax(attn_scores_2, dim=0)
context_vec_2 = attn_weights_2 @ inputs # 形状 (3,)
# -> tensor([0.4419, 0.6515, 0.5683])
\end{lstlisting}
\end{codebox}

\subsection{把 6 个查询者一次算完}

把上面对单个 $\vect{x}^{(2)}$ 做的事，对 6 个词\emph{同时}做。
让 $X$ 表示 $(6, 3)$ 的输入矩阵，则

\[
  \Omega = X X^\top \in \mathbb{R}^{6\times 6}, \qquad
  A = \mathrm{softmax}_{\text{行}}(\Omega), \qquad
  Z = A X \in \mathbb{R}^{6\times 3}.
\]

\begin{codebox}
\begin{lstlisting}
attn_scores  = inputs @ inputs.T           # (6, 6)
attn_weights = torch.softmax(attn_scores, dim=-1)  # 每行 sum=1
all_context_vecs = attn_weights @ inputs   # (6, 3)
\end{lstlisting}
\end{codebox}

\begin{keypointblock}
注意力 = 「拿点积当相似度，softmax 当权重，加权求和当输出」。
本节没有任何可学习参数，只是把这套套路写清楚。
\end{keypointblock}

\section{第二步：把可训练参数加进去（§3.4）}

\subsection{背景问题}

§3.3 的版本有个尴尬：「相似度」完全由原始词向量决定，模型没有任何「调节空间」。
真实任务里我们希望模型自己学到：
「在做这个任务时，词与词之间到底应该看哪些维度才算\emph{有关}？」

\subsection{朴素直觉：把每个词分裂成三个角色}

\begin{intuitionblock}
还是那个开会的比喻。每个人现在被要求扮演\emph{三个角色}：

\begin{itemize}
  \item \textbf{Query（查询者）}：「我想问什么？」
  \item \textbf{Key（钥匙）}：「关于我的信息，可以用什么标签找到我？」
  \item \textbf{Value（值）}：「如果你真的看了我，我会贡献给你什么？」
\end{itemize}

判断「该不该听某人」用 \emph{query 和 key} 做点积。
真正搬过来的内容是 \emph{value}。
而把同一个原始向量分裂成这三个版本的，是三个可训练的矩阵 $W_q, W_k, W_v$。
\end{intuitionblock}

\subsection{数学表达：缩放点积注意力}

设输入维度 $d_{\text{in}}$，输出维度 $d_{\text{out}}$，三个投影矩阵
\[
  W_q, W_k, W_v \in \mathbb{R}^{d_{\text{in}} \times d_{\text{out}}}.
\]
对每个词 $\vect{x}^{(i)}$ 算出三个版本：
\[
  \vect{q}^{(i)} = \vect{x}^{(i)} W_q,\quad
  \vect{k}^{(i)} = \vect{x}^{(i)} W_k,\quad
  \vect{v}^{(i)} = \vect{x}^{(i)} W_v.
\]
注意力分数和权重写成
\[
  \alpha_{ij} \;=\; \mathrm{softmax}_j\!\left(\frac{\vect{q}^{(i)} \cdot \vect{k}^{(j)}}{\sqrt{d_k}}\right),
  \qquad d_k = d_{\text{out}}.
\]
最终
\[
  \vect{z}^{(i)} \;=\; \sum_j \alpha_{ij}\, \vect{v}^{(j)}.
\]
矩阵形式（这就是 transformer 论文那个出名的式子）：
\[
  \boxed{\;
  \mathrm{Attention}(Q, K, V) \;=\; \mathrm{softmax}\!\left(\frac{QK^\top}{\sqrt{d_k}}\right) V.
  \;}
\]

\subsection*{为什么除以 $\sqrt{d_k}$？}

\begin{intuitionblock}
$\vect{q}\cdot\vect{k}$ 是 $d_k$ 项相加。维度 $d_k$ 越大，
点积的数值波动也越大，softmax 容易把几乎全部概率塞给单独一个位置，
变成「非黑即白」的硬选择，梯度会几乎消失。
除以 $\sqrt{d_k}$ 把方差拉回正常范围，softmax 才有「软选择」的余地，模型才训得动。
\end{intuitionblock}

\subsection{代码对应：\texttt{SelfAttention\_v1} 与 \texttt{\_v2}}

notebook 里把上面这套封装成两个 \texttt{nn.Module} 类，结构完全一致，区别只在「投影矩阵怎么写」：

\begin{codebox}
\begin{lstlisting}
class SelfAttention_v1(nn.Module):       # 用裸的 nn.Parameter
    def __init__(self, d_in, d_out):
        super().__init__()
        self.W_query = nn.Parameter(torch.rand(d_in, d_out))
        self.W_key   = nn.Parameter(torch.rand(d_in, d_out))
        self.W_value = nn.Parameter(torch.rand(d_in, d_out))
    def forward(self, x):
        Q, K, V = x @ self.W_query, x @ self.W_key, x @ self.W_value
        scores  = Q @ K.T
        weights = torch.softmax(scores / K.shape[-1] ** 0.5, dim=-1)
        return weights @ V

class SelfAttention_v2(nn.Module):       # 用更标准的 nn.Linear
    def __init__(self, d_in, d_out, qkv_bias=False):
        super().__init__()
        self.W_query = nn.Linear(d_in, d_out, bias=qkv_bias)
        self.W_key   = nn.Linear(d_in, d_out, bias=qkv_bias)
        self.W_value = nn.Linear(d_in, d_out, bias=qkv_bias)
    def forward(self, x):
        Q, K, V = self.W_query(x), self.W_key(x), self.W_value(x)
        scores  = Q @ K.T
        weights = torch.softmax(scores / K.shape[-1] ** 0.5, dim=-1)
        return weights @ V
\end{lstlisting}
\end{codebox}

\begin{remark}
\texttt{v1} 和 \texttt{v2} 输出的数不一样，\emph{不是}因为算法不同，
而是因为 \texttt{nn.Linear} 的默认初始化方式比 \texttt{torch.rand} 好得多。
真实代码里都用 \texttt{v2}。
\end{remark}

\section{第三步：让模型「不许偷看未来」（§3.5）}

\subsection{背景问题}

GPT 是\emph{自回归}的：在预测第 $i$ 个词时，它只能看 $1, 2, \dots, i$ 这些位置，
\textbf{绝对不能看到 $i$ 以后的词}——因为训练时那些词正是它要预测的答案。
但 §3.4 的注意力对 6 个位置一视同仁，没有任何方向感。我们要给它加「方向」。

\subsection{朴素直觉：把右上角抹掉}

\begin{intuitionblock}
注意力权重是一个 $(6, 6)$ 的方阵：第 $i$ 行表示「位置 $i$ 在看哪些位置」。
要让它「不许向右看」，只需保证第 $i$ 行里第 $i+1, i+2, \dots$ 列的权重都是 0。
也就是把整个矩阵的\emph{右上三角}抹成 0，剩下一个下三角。
\end{intuitionblock}

\subsection{两种实现方式}

\textbf{做法 A（教学版）}：先 softmax 再乘下三角 mask，最后\emph{重新归一化}让每行和等于 1。
能跑，但要做两次归一化，啰嗦。

\textbf{做法 B（标准版）}：在 softmax \emph{之前}，把上三角的分数直接设成 $-\infty$。

\begin{mathbox}{为什么 $-\infty$ 这么神奇？}
softmax 第一步是 $\exp(\cdot)$，而 $\exp(-\infty) = 0$。
所以「-∞ 的分数」会自动得到「0 的权重」，剩下的位置自然加起来等于 1，
只要做一次 softmax 就完成了「屏蔽 + 归一化」。
\end{mathbox}

\subsection{代码对应}

\begin{codebox}
\begin{lstlisting}
mask = torch.triu(torch.ones(context_length, context_length), diagonal=1)
# mask 是上三角全 1、对角及以下全 0
masked_scores = scores.masked_fill(mask.bool(), -torch.inf)
weights = torch.softmax(masked_scores / d_k ** 0.5, dim=-1)
\end{lstlisting}
\end{codebox}

\subsection{再加一层 dropout（§3.5.2）}

\begin{intuitionblock}
\textbf{Dropout} 是常见的正则化：训练时随机把一部分数值置 0，让模型不要过度依赖任何单一连接。
这里把 dropout 直接\emph{作用在注意力权重矩阵上}：随机丢弃一些「我在看你」的连线。

为了让被保留下来的位置「期望值不变」，PyTorch 会把保留的元素乘 $\dfrac{1}{1-p}$
（$p$ 是 dropout 比率）。例如 $p = 0.5$ 时，留下来的元素都 $\times 2$。
\end{intuitionblock}

\subsection{封装成 \texttt{CausalAttention}（§3.5.3）}

到这里我们把「可训练 + 因果 + dropout + 支持 batch」打包成一个干净的类：

\begin{codebox}
\begin{lstlisting}
class CausalAttention(nn.Module):
    def __init__(self, d_in, d_out, context_length, dropout, qkv_bias=False):
        super().__init__()
        self.d_out = d_out
        self.W_query = nn.Linear(d_in, d_out, bias=qkv_bias)
        self.W_key   = nn.Linear(d_in, d_out, bias=qkv_bias)
        self.W_value = nn.Linear(d_in, d_out, bias=qkv_bias)
        self.dropout = nn.Dropout(dropout)
        self.register_buffer(
            "mask",
            torch.triu(torch.ones(context_length, context_length), diagonal=1)
        )

    def forward(self, x):                          # x: (b, T, d_in)
        b, T, _ = x.shape
        Q, K, V = self.W_query(x), self.W_key(x), self.W_value(x)
        scores = Q @ K.transpose(1, 2)              # (b, T, T)
        scores.masked_fill_(self.mask.bool()[:T, :T], -torch.inf)
        weights = torch.softmax(scores / K.shape[-1] ** 0.5, dim=-1)
        weights = self.dropout(weights)
        return weights @ V                          # (b, T, d_out)
\end{lstlisting}
\end{codebox}

三个值得注意的小细节：

\begin{itemize}
  \item \texttt{register\_buffer}：mask 不是要训练的参数，但要跟着模型一起搬到 GPU，所以注册为 buffer。
  \item \texttt{transpose(1, 2)}：现在多了一个 batch 维，所以转置的是后两维。
  \item 切片 \texttt{[:T, :T]}：mask 是按最大可能长度建的，实际输入比较短时切一块出来用。
\end{itemize}

\section{第四步：从一个头到多头（§3.6）}

\subsection{背景问题}

一个注意力头只能学一种「关注模式」。但语言里同时存在很多种关系：
词法依赖、指代、修饰、远距离呼应……一个头很难全管。
\textbf{多头注意力}的想法是：让若干个头\emph{并行}地学不同的关注模式，最后把它们的结果拼起来。

\subsection{教学版：把多个 \texttt{CausalAttention} 摞起来（§3.6.1）}

\begin{codebox}
\begin{lstlisting}
class MultiHeadAttentionWrapper(nn.Module):
    def __init__(self, d_in, d_out, context_length, dropout, num_heads, qkv_bias=False):
        super().__init__()
        self.heads = nn.ModuleList([
            CausalAttention(d_in, d_out, context_length, dropout, qkv_bias)
            for _ in range(num_heads)
        ])
    def forward(self, x):
        return torch.cat([h(x) for h in self.heads], dim=-1)
\end{lstlisting}
\end{codebox}

\begin{intuitionblock}
\textbf{它清楚但慢。}\\
$N$ 个头要跑 $N$ 次循环，硬件白白闲着。下面这个生产版会把 $N$ 个头压进\emph{一次}矩阵乘法。
\end{intuitionblock}

\subsection{生产版：权重切分 + 一次矩阵乘（§3.6.2）}

核心想法：

\begin{enumerate}
  \item 还是用一个大的 \texttt{nn.Linear(d\_in, d\_out)} 做 $Q, K, V$，但要求 \texttt{d\_out} 已经能被 \texttt{num\_heads} 整除。
  \item 把最后一维 \texttt{d\_out} \emph{重塑}成 \texttt{(num\_heads, head\_dim)}。
  \item \emph{转置}让 \texttt{num\_heads} 跑到前面，作为新的「批量轴」。
  \item 这样「每个头各自做缩放点积注意力」就变成一次带 batch 的矩阵乘法。
  \item 算完再转置回去、压平最后两维，最后过一个输出投影 \texttt{out\_proj}。
\end{enumerate}

\begin{mathbox}{形状变换}
\[
  \underbrace{(b, T, d_{\text{out}})}_{\text{投影后}}
  \xrightarrow{\text{view}}
  (b, T, h, d_h)
  \xrightarrow{\text{transpose}(1,2)}
  (b, h, T, d_h),
\]
其中 $h$ 是头数，$d_h = d_{\text{out}} / h$。
\end{mathbox}

\begin{codebox}
\begin{lstlisting}
class MultiHeadAttention(nn.Module):
    def __init__(self, d_in, d_out, context_length, dropout, num_heads, qkv_bias=False):
        super().__init__()
        assert d_out % num_heads == 0
        self.d_out, self.num_heads = d_out, num_heads
        self.head_dim = d_out // num_heads
        self.W_query = nn.Linear(d_in, d_out, bias=qkv_bias)
        self.W_key   = nn.Linear(d_in, d_out, bias=qkv_bias)
        self.W_value = nn.Linear(d_in, d_out, bias=qkv_bias)
        self.out_proj = nn.Linear(d_out, d_out)
        self.dropout = nn.Dropout(dropout)
        self.register_buffer(
            "mask",
            torch.triu(torch.ones(context_length, context_length), diagonal=1)
        )

    def forward(self, x):                              # x: (b, T, d_in)
        b, T, _ = x.shape
        Q = self.W_query(x).view(b, T, self.num_heads, self.head_dim).transpose(1, 2)
        K = self.W_key(x)  .view(b, T, self.num_heads, self.head_dim).transpose(1, 2)
        V = self.W_value(x).view(b, T, self.num_heads, self.head_dim).transpose(1, 2)
        # 现在 Q, K, V 形状都是 (b, num_heads, T, head_dim)

        scores = Q @ K.transpose(2, 3)                 # (b, h, T, T)
        scores.masked_fill_(self.mask.bool()[:T, :T], -torch.inf)
        weights = torch.softmax(scores / K.shape[-1] ** 0.5, dim=-1)
        weights = self.dropout(weights)

        ctx = (weights @ V).transpose(1, 2)            # (b, T, h, d_h)
        ctx = ctx.contiguous().view(b, T, self.d_out)  # 拼回 d_out
        return self.out_proj(ctx)                      # (b, T, d_out)
\end{lstlisting}
\end{codebox}

\begin{remark}
注意 \texttt{MultiHeadAttention} 和 \texttt{MultiHeadAttentionWrapper} 的输出维度约定不同：
前者的 \texttt{d\_out} 已经是「所有头加起来」的总维度，
后者是「每个头独立的 \texttt{d\_out}，再拼起来」。所以同一组参数下，
前者输出 \texttt{(2, 6, 2)}，后者输出 \texttt{(2, 6, 4)}。
\end{remark}

\section{把整章压成一张表}

\begin{center}
\renewcommand{\arraystretch}{1.2}
\begin{tabular}{@{}lccccc@{}}
\toprule
小节 & 写法 & 可训练 & 因果 & dropout & 多头 \\
\midrule
§3.3.1 & 手写循环 / 点积                    & 否 & 否 & 否 & 1 \\
§3.3.2 & \texttt{inputs @ inputs.T}        & 否 & 否 & 否 & 1 \\
§3.4.1 & 手写 $W_q, W_k, W_v$              & 是 & 否 & 否 & 1 \\
§3.4.2 & \texttt{SelfAttention\_v1/v2}      & 是 & 否 & 否 & 1 \\
§3.5.1 & + 因果 mask（$-\infty$）          & 是 & 是 & 否 & 1 \\
§3.5.2 & + dropout                          & 是 & 是 & 是 & 1 \\
§3.5.3 & \texttt{CausalAttention}（带 batch）& 是 & 是 & 是 & 1 \\
§3.6.1 & \texttt{MultiHeadAttentionWrapper}  & 是 & 是 & 是 & N（堆叠）\\
§3.6.2 & \texttt{MultiHeadAttention}         & 是 & 是 & 是 & N（并行）\\
\bottomrule
\end{tabular}
\end{center}

\section{这一节真正该记住的}

\begin{enumerate}[label=\textbf{\arabic*.}]
  \item 注意力的本质是\textbf{加权汇总}：
        权重来自 $\mathrm{softmax}(QK^\top / \sqrt{d_k})$，
        汇总的对象是 $V$。
  \item $W_q, W_k, W_v$ 是\textbf{可学习参数}，让模型能为不同任务挑出该看的维度。
  \item $\sqrt{d_k}$ 是\textbf{为了不让 softmax 极化}，否则梯度会消失。
  \item 因果掩码用 $-\infty$ 而不是 $0$，是因为它要在 softmax \emph{之前}生效，
        利用 $\exp(-\infty) = 0$ 自动屏蔽未来。
  \item dropout 用在\textbf{注意力权重矩阵上}，是为了别让模型死记某条「位置 $i \to j$」的连线。
  \item 多头注意力 = 把 $d_{\text{out}}$ 拆成「头数 × 每头维度」，然后让\emph{头}充当批量轴，
        一次矩阵乘把所有头算完，最后再过一次 \texttt{out\_proj} 把它们融合起来。
  \item \texttt{MultiHeadAttention} 这个类，就是第 4 章 GPT 模型里的注意力积木。
\end{enumerate}

\begin{keypointblock}
读完这一章，你应该能在白板上不查代码地写出：\\
1) 缩放点积注意力的公式；\\
2) 因果掩码为什么用 $-\infty$；\\
3) 多头注意力怎么用「reshape + transpose」一次算完。
\end{keypointblock}

\end{document}
